\documentclass[11pt]{article}
\usepackage[hmargin=3cm,bottom=3cm]{geometry}
\usepackage{amsmath}
\usepackage{amssymb}
\usepackage{graphicx,color}
\usepackage[font=scriptsize]{caption}
\usepackage{empheq}
\usepackage{listings}
\usepackage{hyperref}
\usepackage[affil-it]{authblk}
\usepackage{lineno}
\usepackage{bm}
\usepackage{bbold}
\usepackage{listings}
\lstset{breaklines=true,basicstyle=\ttfamily, columns=fullflexible, keepspaces=true}

\setcounter{secnumdepth}{3}

\newcommand*{\defeq}{\mathrel{\vcenter{\baselineskip0.5ex \lineskiplimit0pt
			\hbox{\scriptsize.}\hbox{\scriptsize.}}}%
	=}
\newcommand\p{\ensuremath{\partial}}
\newcommand{\bx}{\mathbf{x}}
\newcommand{\bxi}{\mathbf{\xi}}
\newcommand{\oC}{\overline C} 
\newcommand{\bs}{\mathbf{s}}
\newcommand{\dt}{\Delta t}
\newcommand{\der}{\mathrm{d}}

\newcommand*\pFqskip{8mu}
\catcode`,\active
\newcommand*\pFq{\begingroup
	\catcode`\,\active
	\def ,{\mskip\pFqskip\relax}%
	\dopFq
}
\catcode`\,12
\def\dopFq#1#2#3#4#5{%
	{}_{#1}F_{#2}\biggl[\genfrac..{0pt}{}{#3}{#4};#5\biggr]%
	\endgroup
}


\title{\textit{ParticleTrackingOF} manual}
% \date{}
\author{Tom\'as Aquino}

\begin{document}

\maketitle

\begin{abstract}
This is a brief manual describing the code structure, main processes implemented, and how to compile and use the code on a UNIX system. ParticleTrackingOF is a C++ library for modeling advective--diffusive--reactive transport using Lagrangian particle tracking. It is based in part on OpenFOAM, a free and open source CFD software (https://www.openfoam.com/). It is designed to be used in conjunction with OpenFOAM for mesh generation and numerical determination of the underlying flow field. This offering is not approved or endorsed by OpenCFD Limited, producer and distributor of the OpenFOAM software and owner of the OPENFOAM®  and OpenCFD® trade marks.
\end{abstract}

\section{Compilation and execution}
\label{s::compilation}

The \textit{ParticleTrackingOF} library is written in C++. Executables using the library must be compiled and linked to external dependencies locally using a compiler compatible with C++17  (note that clang++ must be used on macOS due to linking issues with OpenFOAM). The code depends on the following external libraries:
%
\begin{itemize}
\item \emph{OpenFOAM}: CFD software (https://www.openfoam.com/)
\item \emph{boost}: C++ library implementing useful features and extensions (boost.org)
\end{itemize}
These should be obtained and set up according to their respective documentations before using the present library. The executables generated by this library are not OpenFOAM applications, although they rely heavily on the OpenFOAM API (\url{https://www.openfoam.com/documentation/guides/latest/api/}) and are designed to be used with previously-generated OpenFOAM meshes and underlying flow fields.

\subsection{Compiling executables}

The following compilable \lstinline{.cpp} file is provided in the \lstinline{src} directory:
%
\begin{itemize}
\item \lstinline{ParticleTrackingOF.cpp} -- Advective-diffusive-reactive particle tracking.
\end{itemize}
By default, the \lstinline{Makefile} provided for compiling the examples uses the C++ compiler obtained from the OpenFOAM environment variable \lstinline{WM_COMPILER}. Note that the same compiler that was used for OpenFOAM must be used here; failure to do so will lead to dynamic linker errors at runtime. OpenFOAM include directories are set by the variable \lstinline{PTHOF}, which is set to \lstinline{$FOAM_SRC} by default. The paths to provided header files are set in \lstinline{PTH} (\lstinline{../include} by default). These can be changed if necessary, and any additional include paths should be added to the \lstinline{INC} variable. If, for example, you want to use OpenFOAM's third-party boost, add \lstinline{-I/$WM_THIRD_PARTY_DIR} to \lstinline{INC}. Executables must be linked to the OpenFOAM library. The corresponding path is set by the \lstinline{PTHOFLIB} variable, which is set to \lstinline{$FOAM_LIBBIN} by default and may also be changed if necessary. Alternatively, include and library directories can be added to the system \lstinline{PATH} variable as usual. Finally, the target location for compiled executables is set in the \lstinline{PTHBIN} variable (\lstinline{../bin} by default).

Different models can be compiled by changing the \lstinline{using namespace} declaration at the top of \lstinline{main} in the \lstinline{.cpp} file, which should be set to \lstinline{using namespace ptof::model_name}, where \lstinline{name} is the name of an existing model, defined in the corresponding namespace in \lstinline{include/BeadPack/BeadPack_Models.h}. Currently, the available models are:
%
\begin{itemize}
\item \lstinline{advection_2d}
\item \lstinline{advection_diffusion_2d}
\item \lstinline{advection_diffusion_decay_catalytic_2d}
\item \lstinline{advection_diffusion_decay_2d}
\item \lstinline{periodic_cartesian_advection_2d}
\item \lstinline{periodic_cartesian_advection_diffusion_2d}
\item \lstinline{periodic_cartesian_advection_diffusion_decay_catalytic_2d}
\item \lstinline{periodic_cartesian_advection_diffusion_decay_2d}
\item \lstinline{advection_3d}
\item \lstinline{advection_diffusion_3d}
\item \lstinline{advection_diffusion_decay_catalytic_3d}
\item \lstinline{advection_diffusion_decay_3d}
\item \lstinline{periodic_cartesian_advection_3d}
\item \lstinline{periodic_cartesian_advection_diffusion_3d}
\item \lstinline{periodic_cartesian_advection_diffusion_decay_catalytic_3d}
\item \lstinline{periodic_cartesian_advection_diffusion_decay_3d}
\item \lstinline{bcc_cartesian_advection}
\item \lstinline{bcc_cartesian_advection_diffusion}
\item \lstinline{bcc_cartesian_advection_diffusion_decay_catalytic}
\item \lstinline{bcc_cartesian_advection_diffusion_decay}
\item \lstinline{bcc_symmetryplanes_advection}
\item \lstinline{bcc_symmetryplanes_advection_diffusion}
\item \lstinline{bcc_symmetryplanes_advection_diffusion_decay_catalytic}
\item \lstinline{bcc_symmetryplanes_advection_diffusion_decay}
\item \lstinline{advection_diffusion_fpt}
\item \lstinline{periodic_cartesian_advection_diffusion_fpt}
\item \lstinline{bcc_cartesian_advection_diffusion_fpt}
\item \lstinline{bcc_symmetryplanes_advection_diffusion_fpt}
\end{itemize}
These names follow the following conventions:
%
\begin{itemize}
\item Models marked \lstinline{2d} are two-dimensional, and models marked \lstinline{3d} are three-dimensional.
\item Models marked \lstinline{decay_catalytic} include a fluid-solid surface decay reaction, with the solid phase not consumed ($A_F+A_S\to A_S$).
\item Models marked \lstinline{decay} (but not \lstinline{catalytic}) include a fluid-solid surface decay reaction, with the solid phase consumed ($A_F+A_S\to \varnothing$).
\item Models marked \lstinline{advection} include advective transport
\item Models marked \lstinline{diffusion} include diffusive transport
\item Models marked \lstinline{periodic_cartesian} support Cartesian boundary conditions, extracted directly from mesh data for boundaries that specify the periodic type.
\item Models marked \lstinline{bcc} are three-dimensional and exemplify the application of different types of periodic boundary conditions (cartesian or based on symmetry planes) for advective--diffusive and purely--advective transport in a body cubic centered beadpack.
\end{itemize}

Bash scripts to automate the model choice and compilation are provided. The script \lstinline{set_model.sh}, to be run from a command line, should be followed by the model name. For example, running \lstinline{./set_model.sh advection_diffusion} sets the namespace \lstinline{ptof::model_advection_diffusion}. To compile the executables, open a command line, \lstinline{cd} into the \lstinline{src} directory, and run \lstinline{./make.sh}, which takes an argument in the same way. It both sets the model and compiles the corresponding executable. The resulting executable files are placed in the \lstinline{bin} folder and have the base name followed by an underscore and the model name (e.g., \lstinline{ParticleTrackingOF_advection_diffusion_2d}). The script \lstinline{make_all.sh} compiles all available models. Note that running \lstinline{make ParticleTrackingOF} directly produces an executable named \lstinline{ParticleTrackingOF} using the currently set model in \lstinline{bin}, and that running \lstinline{make.sh} or \lstinline{make_all.sh} will delete this file.

If any of the shell scripts need to be given executing permission, this can be achieved as usual by running, e.g.,
%
\begin{lstlisting}[language=bash]
  chmod +x make.sh
\end{lstlisting}

\subsection{Parameters}

In the compiled executables, there are seven types of parameter sets. Default values, if available, are used when \lstinline{`'} is passed for the corresponding parameter. First, the executable itself takes the following parameters (default in $[]$):
%
\begin{itemize}
\item Cases directory $[$../cases$]$
\item Name of case
\item Name of transport parameter set
\item Name of reaction parameter set
\item Name of solver parameter set
\item Name of initial condition parameter set
\item Name of output parameter set
\item Output directory $[\langle$Case directory$\rangle$/output]
\end{itemize}
An executable named \lstinline{ExampleExecutable} is run from a command line by executing \lstinline{./ExampleExecutable} followed by these parameters, separated by spaces. The case name refers to the directory with information about the problem to be solved, which should exist in the cases directory. Running the example executables followed by -h or --help prints a brief description and a list of the required parameters for the executable itself, and for each of the remaining five parameter sets (transport, reaction, solver, initial condition, and output). The required parameters for each of these five parameter sets should be placed in a file in the \lstinline{parameters} directory within the case, named \lstinline{parameters_set_name.dat}, where \lstinline{set} is each of the sets (transport, reaction, etc.) and \lstinline{name} is the corresponding name passed to the executable. In addition, the file \lstinline{of.dat}, within the same directory, must contain:
%
\begin{itemize}
\item OpenFOAM cases directory $[$\lstinline{$FOAM_RUN}$]$
\item OpenFOAM case name $[$\lstinline{$FOAM_DIR}$]$
\item OpenFOAM case time $[$last OpenFOAM case time$]$
\end{itemize}
The OpenFOAM case and time are used to read the previously-generated mesh and flow field to be used in the transport simulation. Values in these files should be placed one per line.

\section{Processes and code structure}

The basic structure of the code relies on a generic CTRW object, which handles a set of Particle objects. Each Particle object is associated with two State objects, describing the particle state associated with the next and previous turning points. The dynamics of these particles are then handled by a Transitions object. Upon a transition, Transition objects take the two particle states, modify the newer state according to the current transition, and set to older state to the previous newer state. For the present use, the Transition object may include any combination of advection, diffusion, and reaction processes, as well as boundary conditions handled by a Boundary object. While the generic CTRW object interface is designed to handle more general dynamics, for the particle tracking random walk (PTRW) dynamics considered here, each transition step is associated with a deterministic time step.

\subsection{Transport processes}

Regarding transport, the BeadPack library focuses on diffusion (with constant diffusion coefficient) and advection according to a heterogeneous flow field. The corresponding particle displacements are handled by a transport-related Transitions object relying on a TimeGenerator object for deterministic time steps. Within this object, particle displacements are handled by a JumpGenerator object, which produces a displacement given the current particle state. JumpGenerator objects implementing various methods for advective and diffusive steps. When both processes are present, a JumpGenerator is used to obtain the full displacement by summing the contributions from the respective JumpGenerators.

For diffusion, the following methods to compute displacement for a given time step are currently implemented:
%
\begin{itemize}
\item Stochastic forward-Euler temporal discretization of the Langevin diffusion equation.
\end{itemize}
All models provided use this option.

For advection, if the flow field is known analytically, methods to use the analytical flow field as a function of position are available in the BeadPack library. However, the flow field must typically be obtained from separate flow simulations. The resulting flow field at an arbitrary set of points in the void space between beads is then to be provided by the user. The BeadPack library allows for computing advective displacements using interpolation to obtain the flow field at arbitrary particle positions. The following interpolation procedures are currently implemented:
%
\begin{itemize}
\item Linear interpolation using a previously-generated OpenFOAM velocity field (based on OpenFOAM routines).
\end{itemize}
For a given method to compute the flow field at an arbitrary location, the following methods to compute advective displacement for a given time step are implemented:
%
\begin{itemize}
\item Forward-Euler temporal discretization
\item Fourth-order Runge--Kutta temporal discretization
\end{itemize}
All models provided use the second option.

\subsection{Reaction processes}

Reaction processes may be added by using a Transitions object which combines a transport-related Transitions object as described above with a reaction, and applies them sequentially. The following reaction processes are currently implemented:
%
\begin{itemize}
\item Fluid-solid reaction on specified boundary patches.
\end{itemize}

The following types of initial distribution of solid reactants are currently implemented:
%
\begin{itemize}
\item \lstinline{uniform}: Homogeneous over the solid phase.
\end{itemize}

\subsection{Boundary conditions}

Boundary conditions are handled by a Boundary object, which takes the previous and current state upon a transport transition, and enforces the boundary conditions on the current state. Boundary conditions are associated with patches in a previously-generated OpenFOAM mesh.
By default, the Boundary object handles the following boundary conditions:
%
\begin{itemize}
\item \lstinline{reflecting}: Reflecting
\item \lstinline{periodic}: Periodic
\item \lstinline{absorbing}: Absorbing
\item \lstinline{info}: Information upon crossing
\item \lstinline{custom}: Enforced by custom Boundary object
\item \lstinline{empty}: No effect (default for unspecified patches in mesh)
\end{itemize}
One of these types must be specified for any patches in the previously-generated OpenFOAM mesh with associated boundary conditions. This information is provided in a file \lstinline{boundary_conditions_dynamics.dat} in the \lstinline{boundary_conditions} directory of a case, where \lstinline{dynamics} stands for different possible choices. Each line of the file should list a mesh patch name followed by one of the boundary condition types. Models not marked \lstinline{fpt} obtain boundary conditions from the file \lstinline{boundary_conditions_transport.dat}. Models marked \lstinline{fpt} obtain them from the file \lstinline{boundary_conditions_firstpassage.dat}. The first should set the appropriate boundary conditions for transport, whereas the second should change the first passage boundaries to \lstinline{absorbing}. The following patch names hold special significance for this library:
%
\begin{itemize}
\item \lstinline{inlet}: Flow inlet
\item \lstinline{outlet}: Flow outlet
\item \lstinline{wallFluidSolid}: Interface of carrier fluid (e.g., water) with solid phase (e.g., grains)
\item \lstinline{wallFluidFluid}: Interface of non-carrier fluid (e.g., air) with solid phase (e.g., grains)
\end{itemize}
The following types are implemented for periodic boundary conditions
%
\begin{itemize}
\item Periodic along cartesian dimensions
\item Periodic according to prescribed symmetry planes
\end{itemize}
The \lstinline{custom} condition has no effect for models not marked \lstinline{fpt}. For models marked \lstinline{fpt}, it reinjects particles according to the initial condition and is meant to be applied to the \lstinline{outlet} patch (assuming it is not periodic). Additional customized boundary conditions can also be set up and implemented. All boundary conditions may also be used to store information in particle states when boundaries are hit or crossed. This is done through Store objects passed to the boundary, and Info objects within the State objects.

\subsection{Initial conditions}

Initial conditions, or more precisely, the injection of particles into the domain, which may happen at any time, are handled by an InitialCondition object.
The following initial conditions are implemented:
%
\begin{itemize}
\item \lstinline{uniform}: Homogeneous throughout the domain
\item \lstinline{flux_weighted}: Flux-weighted throughout the domain
\item \lstinline{uniform_inlet}: Homogeneous at the inlet
\item \lstinline{flux_weighted_inlet}: Flux-weighted at the inlet
\item \lstinline{uniform_solid}: Homogeneous at the solid surface
\item \lstinline{uniform_near_solid}: Homogeneous at a fixed distance to the solid interface
\item \lstinline{uniform_inlet_continuous}: Continuous injection homogeneous at the inlet
\item \lstinline{flux_weighted_inlet_continuous}: Continuous injection flux-weighted at the inlet
\end{itemize}

\subsection{Output}

In the course of particle tracking simulations, quantities of interest relating to particle states can be computed, processed, and output. Typically, these quantities need to be measured a certain number times and/or when certain conditions are met, such as a given time being reached or a given control plane being crossed. The particle state quantities can be accessed for the current and old state of each particle through the CTRW or PTRW objects. By default, the Output class can make measurements with the following options:
%
\begin{itemize}
\item \lstinline{step}: At fixed time intervals, starting at specified time
\item \lstinline{linear}: Fixed number of linearly-spaced measurements, starting and ending at specified times
\item \lstinline{log}: Fixed number of log-spaced measurements, starting and ending at specified times
\end{itemize}
It can also make measurements at the end of the simulation, and it controls the criterion for ending a simulation. The following criteria are implemented:
%
\begin{itemize}
\item \lstinline{time}: Specified time
\item \lstinline{time_max}: Maximum output time
\item \lstinline{mass_below}: Total mass below value
\item \lstinline{mass_above}: Total mass above value
\item \lstinline{all_absorbed}: All particles absorbed
\item \lstinline{one_absorbed}: One particle absorbed
\end{itemize}
It optionally outputs any number of the following:
%
\begin{itemize}
\item \lstinline{position}: Time and particle tags, positions, and masses
\item \lstinline{position_mean}: Time and mean position
\item \lstinline{position_second_moment}: Time and position second moment
\item \lstinline{position_variance}: Time and position variance
\item \lstinline{mass}: Time and total mass
\item \lstinline{velocity}: Time and particle tags and local velocities
\item \lstinline{velocity_gradient}: Time and particle tags and local velocity gradient components
\item \lstinline{pressure}: Time and particle tags and local pressures
\item \lstinline{position_periodic}: Time and particle tags, true positions accounting for periodicity, and masses
\item \lstinline{position_mean_periodic}: Time and true mean position accounting for periodicity
\item \lstinline{position_second_moment_periodic}: Time and true position second moment accounting for periodicity
\item \lstinline{position_variance_periodic}: Time and true position variance accounting for periodicity
\item \lstinline{absorption_time}: Particle absorption times, tags, and masses at end of dynamics
\end{itemize}
Particle tags are unique non-negative integer identifiers. Note that, if desired, the velocity gradient and pressure must be previously computed and present in the OpenFOAM case.

Each of the previous output types is written to a separate file. The output file names have the structure:
%
\begin{lstlisting}
  Data_output_M_model_C_case_OF_ofcase_T_t_R_r_S_s_I_i_O_o_RUN_run.dat
\end{lstlisting}
where \lstinline{output} is the output type, \lstinline{model} is the model name, \lstinline{case} is the simulation case name, \lstinline{ofcase} is the OpenFOAM case name, \lstinline{run} is the run number, and \lstinline{t}, \lstinline{r}, \lstinline{s}, \lstinline{i}, and \lstinline{o} are the parameter set names for transport, reaction, solvers, initial condition, and output, respectively.

%\bibliography{ParticleTrackingOF}
%\bibliographystyle{unsrt}

\end{document}